\section{Mina erfarenheter från LIA-perioden}
\label{sec:erfarenheter}

Under LIA-perioden fick jag många nya erfarenheter, både tekniska och personliga. En av de mest betydelsefulla erfarenheterna var när jag arbetade med att deploya Dagster och Postgres till Kubernetes. Själva deploy-processen gick ganska smidigt, men eftersom detta var första gången jag arbetade med Dagster så var det en utmaning att förstå hur systemet fungerade. Det var särskilt svårt i början när jag stötte på fel och behövde felsöka, eftersom jag ännu inte hade någon tidigare erfarenhet av Dagster.

Med tiden lärde jag mig dock hur jag skulle läsa felmeddelanden, testa olika lösningar och felsöka på ett strukturerat sätt. Jag märkte att jag blev lugnare när problem uppstod och började se troubleshooting som en naturlig del av arbetet. Denna utveckling gav mig ökat självförtroende, eftersom jag kände att jag faktiskt kunde lösa riktiga problem i en professionell miljö.

Min första ticket var också en viktig del av mina erfarenheter. Det var ett problem som Insightas Dagster-team hade definierat, och det var min uppgift att utvärdera om K8sRunLauncher kunde användas för att deploya Dagster utan att avbryta pågående jobb. Det var ett stort ansvar, men jag lyckades genomföra hela uppgiften och uppfylla alla acceptance criteria. Det kändes väldigt bra att teamet litade på mig och att jag kunde leverera något som var användbart för dem.

Om jag skulle göra samma uppgift igen skulle jag börja med att lära mig ännu lite mer om hur Dagster kommunicerar med Kubernetes. Jag förstod senare att om man redan kan grunderna i det samspelet så blir felsökningen mycket enklare. Även om Dagster inte är ett typiskt DevOps-verktyg så insåg jag att man som DevOps-ingenjör ändå behöver förstå plattformen eftersom den är en central del av dataplattformen som hela företaget använder.

Sammanfattningsvis har mina erfarenheter från LIA-perioden lärt mig att kombinera nya tekniska kunskaper med problemlösning i verkliga projekt. Jag har fått känna hur det är att jobba på riktigt, där man måste hitta lösningar, samarbeta med teamet och ta ansvar för sina uppgifter.